\Needspace{11\baselineskip}
\begin{center}
\textit{Mindestens der Hälfte aller für eine gegebene Determinante zuweisbaren Charaktere können keine eigentlich primitiven Geschlechter entsprechen (bei negativer Determinante keine positiven).}
\end{center}

\Needspace{8\baselineskip}
\Article{261}
\textsc{Satz.} \textit{Bei einer positiven Nichtquadratzahl als Determinante können mindestens der Hälfte aller zuweisbaren Charaktere keine eigentlich primitiven Geschlechter entsprechen; bei negativer Determinante keine positiven eigentlich primitiven Geschlechter.}

\textit{Beweis.} Es sei $m$ die Anzahl aller eigentlich primitiven (positiven) Geschlechter der Determinante $D$ und $k$ die Anzahl der Klassen in jedem Geschlecht, sodass $km$ die Anzahl aller eigentlich primitiven (positiven) Klassen ist; $n$ sei die Anzahl aller verschiedenen für diese Determinante zuweisbaren Charaktere. Dann ist
\SourceLeaf{297}
nach Artikel 258 die Anzahl aller eigentlich primitiven (positiven) ambigen Klassen gleich $\frac12 n$. Somit ist nach dem vorigen Artikel die Anzahl aller eigentlich primitiven Klassen, die durch Verdopplung einer solchen Klasse entstehen können, gleich $\frac{2km}{n}$. Nach Artikel 247 gehören aber alle diese Klassen zum Hauptgeschlecht, das $k$ Klassen enthält. Können also alle Klassen des Hauptgeschlechts durch Verdopplung irgendeiner Klasse entstehen (was, wie wir im Folgenden beweisen werden, tatsächlich immer gilt), so ist $\frac{2km}{n}=k$ oder $m=\frac12 n$; gewiss kann aber weder $\frac{2km}{n}>k$ noch daher $m>\frac12 n$ sein. Da somit die Anzahl aller eigentlich primitiven (positiven) Geschlechter gewiss nicht größer als die Hälfte aller zuweisbaren Charaktere ist, können mindestens der Hälfte dieser Charaktere keine solchen Geschlechter entsprechen. \textit{Damit ist der Satz bewiesen.} --- Dabei muss man beachten, dass daraus noch nicht folgt, dass der Hälfte aller zuweisbaren Charaktere tatsächlich eigentlich primitive (positive) Geschlechter entsprechen. Die Wahrheit dieser sehr wichtigen Aussage lässt sich erst weiter unten aus den tiefsten Geheimnissen der Zahlen erschließen.

Da es bei negativer Determinante stets ebenso viele negative wie positive Geschlechter gibt, können offenbar von allen zuweisbaren Charakteren nicht mehr als die Hälfte zu negativen eigentlich primitiven Geschlechtern gehören. Darüber und über uneigentlich primitive Geschlechter werden wir weiter unten sprechen. Schließlich gilt der Satz nicht für positive Quadratzahlen als Determinanten; bei diesen sieht man ohne Mühe, dass jedem zuweisbaren Charakter tatsächlich Geschlechter entsprechen.

\Needspace{11\baselineskip}
\begin{center}
\textit{Ein zweiter Beweis des Fundamentalsatzes und der übrigen Sätze über die Reste $-1$, $+2$, $-2$.}
\end{center}

\Needspace{8\baselineskip}
\Article{262}
Lassen sich also für eine gegebene Nichtquadratzahl $D$ als Determinante nur zwei verschiedene Charaktere zuweisen, so entspricht nur einem ein eigentlich primitives (positives) Geschlecht (das nur das Hauptgeschlecht sein kann); der andere gehört zu keiner eigentlich primitiven (positiven) Form dieser Determinante. Das geschieht bei den Determinanten $-1$, $2$, $-2$, $-4$. den positiv genommenen Primzahlen der Form $4n+1$ und den negativ genommenen der Form $4n+3$, schließlich bei allen positiv genommenen Potenzen ungeraden Exponenten von Primzahlen der Form $4n+1$ und bei den Potenzen von Primzahlen der Form $4n+3$, die je nach geradem oder ungeradem Exponenten positiv oder negativ genommen werden. Aus diesem Prinzip können wir eine neue Methode gewinnen, nicht nur den Fundamentalsatz, sondern auch die übrigen Sätze des vorigen Abschnitts über die Reste $-1$, $+2$, $-2$ zu beweisen. Sie ist von den im
\SourceLeaf{298}
vorigen Abschnitt benutzten Methoden völlig verschieden und scheint ihnen an Eleganz keineswegs nachzustehen. Die Determinante $-4$ und diejenigen, die Primzahlpotenzen sind, lassen wir jedoch aus, da sie nichts Neues lehren.

Für die Determinante $-1$ gibt es also keine positive Form mit dem Charakter $3,4$; für die Determinante $+2$ überhaupt keine Form mit dem Charakter $3\ \mathit{und}\ 5,8$; bei der Determinante $-2$ hat keine positive Form den Charakter $5\ \mathit{und}\ 7,8$; bei der Determinante $+p$, wenn $p$ eine Primzahl der Form $4n+1$ ist, oder bei der Determinante $-p$, wenn $p$ eine Primzahl der Form $4n+3$ ist, hat keine eigentlich primitive Form (im zweiten Fall keine positive) den Charakter $Np$. Damit beweisen wir die Sätze des vorigen Abschnitts folgendermaßen.

I. $-1$ ist Nichtrest jeder (positiven) Zahl der Form $4n+3$. Wäre nämlich $-1$ Rest einer solchen Zahl $A$, so wäre bei $-1=BB-AC$ die Form $(A,B,C)$ eine positive Form der Determinante $-1$ mit dem Charakter $3,4$.

II. $-1$ ist Rest jeder Primzahl $p$ der Form $4n+1$. Denn der Charakter der Form $(-1,0,p)$ ist, wie bei allen eigentlich primitiven Formen der Determinante $p$. gleich $Rp$, also $-1Rp$.

III. Sowohl $+2$ als auch $-2$ ist Rest jeder Primzahl $p$ der Form $8n+1$. Denn entweder sind die Formen $\left(8,\allowbreak 1,\allowbreak \frac{1-p}{8}\right)$, $\left(-8,\allowbreak 1,\allowbreak \frac{p-1}{8}\right)$ oder die Formen $\left(8,\allowbreak 3,\allowbreak \frac{9-p}{8}\right)$, $\left(-8,\allowbreak 3,\allowbreak \frac{p-9}{8}\right)$ eigentlich primitiv (je nachdem, ob $n$ ungerade oder gerade ist); ihr Charakter ist also $Rp$. Daher gilt $+8Rp$ und $-8Rp$, woraus auch $2Rp$, $-2Rp$ folgt.

IV. $+2$ ist Nichtrest jeder Zahl der Form $8n+3$ oder $8n+5$. Wäre es nämlich Rest einer solchen Zahl $A$, so gäbe es eine Form $(A,B,C)$ der Determinante $+2$ mit dem Charakter $3\ \mathit{und}\ 5,8$.

V. Ebenso ist $-2$ Nichtrest jeder Zahl der Form $8n+5$ oder $8n+7$; sonst gäbe es eine Form $(A,B,C)$ der Determinante $-2$ mit dem Charakter $5\ \mathit{und}\ 7,8$.

VI. $-2$ ist Rest jeder Primzahl $p$ der Form $8n+3$. Diese Aussage lässt sich auf zwei Arten beweisen. \textit{Erstens}: Da nach IV $+2Np$ und nach I $-1Np$ gilt, muss $-2Rp$ gelten. Der \textit{zweite} Beweis ergibt sich aus der Betrachtung der Determinante $+2p$, für die vier Charaktere zuweisbar sind, nämlich $Rp$, $1\ \mathit{und}\ 3,8$; $Rp$, $5\ \mathit{und}\ 7,8$; $Np$, $1\ \mathit{und}\ 3,8$; $Np$, $5\ \mathit{und}\ 7,8$. Mindestens zweien davon entsprechen also keine Geschlechter. Nun hat die Form
\SourceLeaf{299}
$(1,0,-2p)$ den ersten Charakter, die Form $(-1,0,2p)$ den vierten; somit müssen der zweite und dritte ausgeschlossen werden. Da also der Charakter der Form $(p,0,-2)$ in Bezug auf die Zahl $8$ gleich $1\ \mathit{und}\ 3,8$ ist, kann ihr Charakter in Bezug auf $p$ nur $Rp$ sein, also $-2Rp$.

VII. $+2$ ist Rest jeder Primzahl $p$ der Form $8n+7$, was sich auf zwei Arten beweisen lässt. \textit{Erstens}: Da aus I und V $-1Np$, $-2Np$ folgt, gilt $+2Rp$. \textit{Zweitens}: Da entweder $\left(8,\allowbreak 1,\allowbreak \frac{1+p}{8}\right)$ oder $\left(8,\allowbreak 3,\allowbreak \frac{9+p}{8}\right)$ eine eigentlich primitive Form der Determinante $-p$ ist (je nachdem, ob $n$ gerade oder ungerade ist), ist ihr Charakter $Rp$, also $8Rp$ und $2Rp$.

VIII. Jede Primzahl $p$ der Form $4n+1$ ist Nichtrest jeder ungeraden Zahl $q$, die Nichtrest von $p$ ist. Wäre nämlich $p$ Rest von $q$, so gäbe es offenbar eine eigentlich primitive Form der Determinante $p$ mit dem Charakter $Np$.

IX. Ist ebenso irgendeine ungerade Zahl $q$ Nichtrest einer Primzahl $p$ der Form $4n+3$, so ist $-p$ Nichtrest von $q$; sonst gäbe es eine positive eigentlich primitive Form der Determinante $-p$ mit dem Charakter $Np$.

X. Jede Primzahl $p$ der Form $4n+1$ ist Rest jeder anderen Primzahl $q$, die Rest von $p$ ist. Ist auch $q$ von der Form $4n+1$, so folgt dies sofort aus VIII; ist dagegen $q$ von der Form $4n+3$, so ist auch $-q$ Rest von $p$ (nach II), also $pRq$ (nach IX).

XI. Ist irgendeine Primzahl $q$ Rest einer anderen Primzahl $p$ der Form $4n+3$, so ist $-p$ Rest von $q$. Ist nämlich $q$ von der Form $4n+1$; so folgt aus VIII $pRq$, also nach II $-pRq$; der Fall aber, dass auch $q$ von der Form $4n+3$ ist, entzieht sich dieser Methode, lässt sich jedoch durch Betrachtung der Determinante $+pq$ leicht erledigen. Von den vier für diese Determinante zuweisbaren Charakteren $Rp$. $Rq$; $Rp$, $Nq$; $Np$, $Rq$; $Np$, $Nq$ können nämlich zweien keine Geschlechter entsprechen. Da die Charaktere der Formen $(1,0,-pq)$, $(-1,0,pq)$ jeweils der erste und vierte sind, können der zweite und dritte zu keiner eigentlich primitiven Form der Determinante $pq$ gehören. Da somit der Charakter der Form $(q,0,-p)$ in Bezug auf die Zahl $p$ nach Voraussetzung $Rp$ ist, muss der Charakter derselben Form in Bezug auf die Zahl $q$ gleich $Rq$ sein, also $-pRq$. \textit{Damit ist der Satz bewiesen.}

Nimmt man in den Aussagen VIII und IX an, dass $q$ eine Primzahl bezeichnet, so ergeben diese zusammen mit X und XI den Fundamentalsatz des vorigen Abschnitts.
\Needspace{18\baselineskip}
\SourceLeaf{300}
\begin{center}
\textit{Die Hälfte der Charaktere, denen keine Geschlechter entsprechen können, wird genauer bestimmt.}
\end{center}

\Needspace{8\baselineskip}
\Article{263}
Nachdem wir den Fundamentalsatz mit einem neuen Beweis bestätigt haben, werden wir zeigen, wie man für jede gegebene Nichtquadratzahl als Determinante jene Hälfte der Charaktere erkennt, denen keine eigentlich primitiven (positiven) Formen entsprechen können. Dies lässt sich umso kürzer erledigen, als die Grundlage bereits in der Untersuchung der Artikel 147--150 enthalten ist. Es sei $ee$ das größte Quadrat, das die vorgelegte Determinante $D$ teilt, und $D=D'ee$, sodass $D'$ keinen quadratischen Faktor enthält. Ferner seien $a$, $b$, $c$ usw. alle ungeraden Primteiler von $D'$; dann ist $D'$ ohne Rücksicht auf das Vorzeichen entweder ihr Produkt oder das Doppelte dieses Produkts. Mit $\Omega$ werde die Menge der Einzelcharaktere $Na$, $Nb$, $Nc$ usw. allein bezeichnet, wenn $D'\equiv1\ (\mathrm{mod.}\ 4)$; unter Hinzunahme des Charakters $3,4$, wenn $D'\equiv3$ und $e$ ungerade oder das Doppelte einer ungeraden Zahl ist; unter Hinzunahme von $3,8$ und $7,8$, wenn $D'\equiv3$ und $e$ durch vier teilbar ist; unter Hinzunahme entweder des Charakters $3\ \mathit{und}\ 5,8$ oder der beiden $3,8$ und $5,8$, wenn $D'\equiv2\ (\mathrm{mod.}\ 8)$, je nachdem, ob $e$ ungerade oder gerade ist; schließlich unter Hinzunahme entweder des Charakters $5\ \mathit{und}\ 7,8$ oder der beiden $5,8$ und $7,8$, wenn $D'\equiv6\ (\mathrm{mod.}\ 8)$, je nachdem, ob $e$ ungerade oder gerade ist. Dann können allen Gesamtcharakteren, die eine ungerade Anzahl von Einzelcharakteren aus $\Omega$ enthalten, keine eigentlich primitiven (positiven) Geschlechter der Determinante $D$ entsprechen. In allen Fällen tragen die Einzelcharaktere, die die Beziehung zu solchen Primteilern von $D$ ausdrücken, die $D'$ nicht teilen, nichts zur Möglichkeit oder Unmöglichkeit von Geschlechtern bei. --- Nach der Kombinationslehre erkennt man sehr leicht, dass so tatsächlich die Hälfte aller zuweisbaren Gesamtcharaktere ausgeschlossen wird.

Diese Regeln werden folgendermaßen bewiesen. Aus den Prinzipien des vorigen Abschnitts oder den im vorigen Artikel erneut bewiesenen Sätzen folgt ohne Mühe: Ist $p$ eine Primzahl (ungerade und positiv), die $D$ nicht teilt und einen der ausgeschlossenen Charaktere hat, so enthält $D'$ eine ungerade Anzahl von Faktoren, die Nichtreste von $p$ sind; somit ist $D'$ und daher auch $D$ Nichtrest von $p$. Ferner kann offenbar ein Produkt aus beliebig vielen zu $D$ teilerfremden ungeraden Zahlen, von denen keine einen ausgeschlossenen Charakter hat, ebenfalls nicht mit einem solchen Charakter übereinstimmen. Umgekehrt enthält daher jede zu $D$ teilerfremde positive ungerade Zahl, die einen der \SourceLeaf{301}ausgeschlossenen Charaktere hat, gewiss einen Primfaktor derselben Art, sodass $D$ Nichtrest dieser Zahl ist. Gäbe es also eine eigentlich primitive (positive) Form der Determinante $D$, die einem ausgeschlossenen Charakter entspricht, so wäre $D$ Nichtrest jeder durch eine solche Form darstellbaren positiven ungeraden und zu ihm teilerfremden Zahl; dies ist offenbar mit dem Satz von Artikel 154 nicht vereinbar.

Als Beispiele vergleiche man die Klassifikationen in den Artikeln 231, 232, deren Anzahl jeder nach Belieben vermehren kann.
